\scanpage{33}
de $\xi'$, donc

\[
\left\{\begin{aligned}
&\lambda\ \text{un autom de }\underline{\mathrm{End}}(G),\\
&v\qquad\rule{2.8em}{.35pt}\qquad E\\
&w\qquad\rule{3em}{.35pt}\qquad G\otimes E
\end{aligned}\right.
\]

avec commutativité de

\diagram{\begin{tikzcd}[ampersand replacement=\&]
\underline{\mathrm{End}}(G)\otimes\underline{\mathrm{End}}(E)
 \arrow[r,"\varphi_{G,E}","\sim"']\arrow[d,"\lambda\otimes\mathrm{End}(v)"']
 \&\underline{\mathrm{End}}(G\otimes E)\arrow[d,"\mathrm{End}(w)"]\\
\underline{\mathrm{End}}(G)\otimes\underline{\mathrm{End}}(E)
 \arrow[r,"\varphi_{G,E}"]
 \&\underline{\mathrm{End}}(G\otimes E)
\end{tikzcd}}

\deleted{On sait} Il suffit de voir que \uline{loc} il provient d'un \deleted{$(u,v)$}
autom de $\xi$, qui sera de la forme $(u,v)$, $u$ un autom
de $G$, donc on veut prouver

\[
\lambda\simeq\mathrm{End}(u),\quad w=u\otimes v
\]

Or on sait que loc. $\lambda$ est \deleted{\unsure{induit}} intérieur, i.e. \unsure{que|qu'|qu'il} \inserted{loc}
\unsure{on} peut trouver $u$ tel $\lambda=\mathrm{End}(u)$. Quitte à multiplier \inserted{$\sigma$} par
$F(u,v)^{-1}$, on peut supposer que $\lambda=\mathrm{id}$. Donc on \deleted{est} a
$\varphi\circ(\mathrm{id}\otimes\mathrm{End}(v))\simeq\mathrm{End}(w)\circ\varphi$, et on doit prouver qu'il existe [loc.]
une section $u$ de $\underline O_{\unsure{T|\mathscr T}}^{*}$ telle que $w=\unsure{u}( \mathrm{id}_{G}\otimes v)$, \deleted{\unsure{or}} \inserted{i.e. $w=\mathrm{id}_{G}\otimes uv$}

\deleted{prenant la sous-algèbre $1_G\otimes\underline{\mathrm{End}}(E)$ de $\underline{\mathrm{End}}(G)\otimes\underline{\mathrm{End}}(\unsure{G|E})$
la commutativité \deleted{\unsure{postule|postulée|partielle}} donne}

\[
\begin{aligned}
&\deleted{\mathrm{End}(w)(\mathrm{id}_G\otimes X)=\mathrm{id}_G\otimes\mathrm{End}(v)(X)\qquad \text{pour }X\in\Gamma\,\mathrm{End}(E)}\\
&\deleted{w(\mathrm{id}_G\otimes X)w^{-1}=\mathrm{id}_G\otimes vXv^{-1}}
\end{aligned}
\]

L'hyp. de compatibilité
s'écrit que pour tt end. $X$ de $G$, $Y$ de $E$, on a

\[
\deleted{\unsure{\varphi(X\otimes Y)}}\quad
w\circ(\underbrace{X\otimes Y}_{\substack{\text{endom}\\\text{produit}\\\text{tensoriel}\\\text{de }G\ \mathrm{et}\ E}})w^{-1}
=\underbrace{X\otimes{}}_{\mathrm{id}}vYv^{-1}
\]

donc $\mathrm{End}(w)$ est un autom. \inserted{connu} de $\mathrm{End}(G\otimes E)$ \inserted{$\simeq\underline{\mathrm{End}}(G)\otimes\underline{\mathrm{End}}(E)$}
\deleted{qui \unsure{commute} à l'action de $\mathrm{End}(G)$, donc elle est déterminée par
son action sur $1_G\otimes\mathrm{End}(E)$ \unsure{donnée}} savoir il
est égal à la troisième composante de $\deleted{F(\mathrm{id}_G,v)}\inserted{\mathrm{End}(\mathrm{id}_G\otimes v)}$, \deleted{Donc}
\deleted{$w$} Donc $w$ en diffère par une \deleted{unité} section, cqfd.

\uline{Corollaire} \deleted{Le lemme reste valable si on remplace $C'$ par la catégorie \unsure{des quadruples} tels}
\inserted{Soit $A$ une Algèbre \deleted{d'Azumaya} sur $\mathscr T$}
\deleted{[loc libre de t.f. \unsure{rg|sg} partout non nulle, $A^{o}\otimes A\xrightarrow{\sim}\underline{\mathrm{End}}(A)$]}
telle qu'il existe $E,F$ modules loc. libres de t.f. \deleted{tel que}\inserted{et un isom}
$A\otimes\underline{\mathrm{End}}(E)\xrightarrow[\sim]{\varphi}\underline{\mathrm{End}}(F)$ \inserted{donné}
Alors \deleted{$A$ est loc} $\exists G$ loc libre de t.f. [partout non nul]
\deleted{tel que}\inserted{et} $A\xrightarrow[\sim]{\psi}\mathrm{End}(G)$, de plus $G,\psi$ est déterminé
à isom unique près.

En effet, l'hyp implique que $A$ est loc libre de t.f. comme
Module, de plus $A^o\otimes A\otimes\mathrm{End}(E)\simeq A^o\otimes\mathrm{End}(F)$
% archive sheet number: 32
